\section{Compléments sur l'apprentissage du réseau bayésien (Apprendre)}
\label{sec:S4}

\subsection{Construction de la base d'apprentissage}
\label{sec:S4.1}

Chaque version d'une crise, c'est-à-dire une crise jouée sous un comportement de gestion donné, fournit une observation ; la base principale en compte 8 000. Les variables sont organisées selon les causes qu'un diagnostic doit pouvoir distinguer : l'exposition du réseau, la perturbation, la politique de gestion et la réponse effective du gestionnaire, complétées par le mécanisme de la chute et par la résilience obtenue (tableau~5 du manuscrit).

La demande de référence n'est pas une variable de contexte : calibrée sur le volume livrable pendant l'incident, elle dépend de la perturbation elle-même. Elle est remplacée par la position de la cible dans la topologie : un site ou une liaison constitue un point de passage obligé si sa suppression déconnecte toutes les sources du client. Sur le réseau ISOMORPH, le hub de Nashville, l'entrepôt d'Atlanta et la liaison qui les relie sont dans ce cas ; la congestion et le pic de demande touchent le réseau entier. Cette variable ne dépend que de la structure du réseau et se transpose à tout réseau étudié. La politique est représentée par le profil seul : ses composantes (seuil, délai, rythme, ordre des leviers) lui sont parfaitement redondantes dans une base où elles sont fixées par profil. Une réponse absente, lorsque le taux de service n'est jamais passé sous le seuil critique, est codée par un état explicite « sans objet ».

Les paramètres de la courbe sont conservés comme variables intermédiaires : ils décrivent le mécanisme par lequel une cause produit son effet (tableau~\ref{tab:S3}). Trois sont écartés faute de variation dans la base : $k$, proche de 1 par construction, $q$, proche de 0 parce que les crises se résorbent avant la fin de la période observée, et $g$, presque toujours nul. Parmi les indicateurs, IPC, fidèle au service perdu et sans ambiguïté de sens, sert de cible au pronostic ; TRP et CRD le complètent. SRA, redondant avec les paramètres, ERR, trop peu dispersé, et PR, dont le sens favorable dépend de la profondeur de la chute, sont écartés.

Toutes les versions sont conservées, y compris celles dont l'ajustement est imparfait : écarter les crises mal ajustées éliminerait surtout les crises courtes, donc une partie des réussites du profil réactif. Les variables continues sont discrétisées en trois classes de même effectif, dont les seuils sont calculés sur l'échantillon d'apprentissage puis figés pour les crises nouvelles.

\subsection{Apprentissage sous contraintes}
\label{sec:S4.2}

L'apprentissage est réalisé avec l'outil BayesArtist \citep{bayesartist}. La structure est apprise par recherche locale gloutonne au critère d'information bayésien (BIC) \citep{schwarz1978}. L'ordre des paliers est imposé comme une liste d'arcs interdits, avec au plus quatre parents par variable, de sorte qu'aucune structure apprise ne peut violer la temporalité de la crise ; cette combinaison de connaissance du domaine et de données suit le principe de \citet{heckerman1995}. Les tables de probabilités sont estimées par comptage avec un lissage de Laplace, qui évite d'attribuer une probabilité nulle à une combinaison de causes absente de la base, et l'inférence est exacte, par élimination de variables ou arbre de jonction \citep{koller2009}.

La politique, intervention comparée sur des crises strictement identiques, a un effet sur la résilience identifiable sans ajustement. Les leviers effectivement mobilisés dépendent au contraire de ce que le gestionnaire a perçu : ce sont des médiateurs.

BayesArtist n'étant pas diffusé, le même apprentissage peut être conduit avec tout logiciel de réseaux bayésiens qui apprend la structure et les paramètres à partir d'un fichier CSV et accepte des arcs interdits, par exemple Hugin, GeNIe, le paquet R bnlearn ou la bibliothèque Python pgmpy ; les algorithmes de recherche différant d'un outil à l'autre, de légers écarts de structure sont possibles.

\subsection{Inférence, préconisation et protocole de validation}
\label{sec:S4.3}

En pronostic, le réseau donne la distribution de la résilience pour un contexte et une politique, par exemple la probabilité d'une perte cumulée faible face à la fermeture sévère d'un site de stockage gérée par un profil prudent. En diagnostic, il remonte d'une résilience dégradée vers les causes les plus probables : point de passage obligé touché, choc difficilement surmontable, perception tardive, inertie décisionnelle ou leviers inadaptés.

La préconisation combine les deux. Face à une nouvelle crise, décrite en tout ou partie, et à un objectif de résilience, chaque politique candidate est évaluée par la probabilité d'atteindre cet objectif, qui se lit directement dans le réseau puisque la politique est une intervention. La règle retenue privilégie la sobriété : parmi les politiques dont cette probabilité dépasse un seuil $\theta$ fixé par le décideur, est préconisée la moins interventionniste, les politiques étant ordonnées par le nombre moyen de leviers qu'elles engagent (sans réaction, tardif, prudent, réactif). Si aucune n'atteint le seuil, la politique la plus probable est préconisée. Cette règle traduit le fait que tout levier a un coût que le modèle ne valorise pas : à résilience égale, mobiliser moins de leviers est préférable.

La même question peut être posée à rebours : fixer l'objectif, par exemple une perte cumulée faible, et le contexte connu, puis lire la distribution du profil. En général, remonter d'un effet vers une décision ne livre qu'une association, car les décisions observées dépendent de la situation. Ici, le profil est une variable racine, assignée selon une loi uniforme indépendamment du contexte : sa probabilité a posteriori est donc proportionnelle à sa probabilité d'atteindre l'objectif, et l'inférence à rebours classe les politiques exactement comme le pronostic. Elle fournit en une requête par contexte une carte de préconisation, qui désigne la politique la plus efficace ; la règle de sobriété reste nécessaire pour retenir la plus économe parmi celles qui suffisent. Cette propriété ne vaut pas pour les leviers, qui dépendent de ce que le gestionnaire perçoit.

La validation exploite une propriété rare de la base : chaque crise y est observée sous les quatre politiques. Les 2 000 scénarios sont partagés aléatoirement en 1 600 scénarios d'apprentissage et 400 de test ; le découpage porte sur les scénarios, les quatre versions d'une même crise n'étant pas indépendantes. Le pronostic est jugé sur la prédiction de la classe de perte cumulée à partir des seules variables connues au moment de décider, par le taux de classement correct et le score de Brier, comparés à la classe majoritaire et à un modèle fondé sur la seule politique. Le diagnostic est jugé en confrontant, pour les crises à forte perte, la distribution a posteriori des causes à leur fréquence observée. La préconisation est jugée directement : chaque crise de test ayant été jouée sous toutes les politiques, on vérifie si la politique préconisée a effectivement atteint l'objectif, et l'on mesure la part des préconisations plus sobres que le profil réactif. Une courbe d'apprentissage, de 200 à 1 600 scénarios, indique si la taille de la base est suffisante.
